% Clean-room contest paper baseline (self-written; no upstream template code is vendored).
% Compile with XeLaTeX from the repository root:
%   xelatex -interaction=nonstopmode -output-directory=paper paper/main.tex   (twice)
\documentclass[12pt,a4paper]{ctexart}
\usepackage[margin=2.5cm]{geometry}
\usepackage{amsmath,amssymb,bm}
\usepackage{graphicx}
\usepackage{booktabs,longtable,tabularx,multirow,array}
\usepackage{float}
\usepackage{caption,subcaption}
\usepackage{enumitem}
\usepackage{xcolor}
\usepackage{listings}
\usepackage[hidelinks]{hyperref}   % 正文无彩色：链接与引用均为黑色
\graphicspath{{paper/figures/}}

\ctexset{
  section/format=\large\bfseries\heiti,
  section/number=\chinese{section},
  subsection/format=\normalsize\bfseries\heiti,
  subsubsection/format=\normalsize\bfseries\heiti,
}

% ---------------------------------------------------------------------------
% 附录代码清单的排版：整幅细框、行号、按语法着色（灰阶，不喧宾夺主）、
% 自动折行并按缩进对齐，中文注释经 ctex 正常显示。
% ---------------------------------------------------------------------------
\lstdefinestyle{pycode}{
  language=Python,
  basicstyle=\ttfamily\scriptsize,
  keywordstyle=\bfseries\color{black},
  commentstyle=\itshape\color{black!55},
  stringstyle=\color{black!70},
  breaklines=true,
  breakatwhitespace=false,
  postbreak=\mbox{$\hookrightarrow$\space},   % 折行续行标记（不能用 \textcolor，会破坏 listings 的间距计算）
  columns=flexible,
  showstringspaces=false,
  keepspaces=true,
  frame=single,
  framexleftmargin=14pt,
  rulecolor=\color{black!45},
  numbers=left,
  numberstyle=\tiny\color{black!50},
  numbersep=7pt,
  extendedchars=true,
  tabsize=4,
  xleftmargin=20pt,
  xrightmargin=2pt,
  lineskip=1.2pt,
  captionpos=b,
}
\lstset{style=pycode}

% Frozen-number macros (injected by scripts/latex_assembly.py)
\newcommand{\qoneCcenteroneeightzerozeros}{2.55\,kg/kg}
\newcommand{\qoneCsurfaceoneeightzerozeros}{1.5102\,kg/kg}
\newcommand{\qoneTcenteroneeightzerozeros}{33.5753\,°C}
\newcommand{\qoneTsurfaceoneeightzerozeros}{36.7856\,°C}
\newcommand{\qoneanalyticalmaxabserr}{$6.18\times10^{-7}$\,kg/kg}
\newcommand{\qonegridrefinementdiff}{$3.24\times10^{-6}$\,kg/kg}
\newcommand{\qonehmplusfivepctsurfaceshift}{-0.0352\,kg/kg}
\newcommand{\qoneindependentimplementationsdiff}{$9.72\times10^{-7}$\,kg/kg}
\newcommand{\qonemassbalancerel}{$5.41\times10^{-11}$}
\newcommand{\qonetimeaccuracytolerancediff}{$3.89\times10^{-10}$\,kg/kg}
\newcommand{\qtwoCcenterthreeh}{1.7662\,kg/kg}
\newcommand{\qtwoCsurfacethreeh}{1.0081\,kg/kg}
\newcommand{\qtwoTcenterthreeh}{49.8495\,°C}
\newcommand{\qtwoTsurfacethreeh}{49.9664\,°C}
\newcommand{\qtwoanalyticalmaxabserr}{$5.98\times10^{-7}$\,kg/kg}
\newcommand{\qtwodryingtime}{57.4647\,h}
\newcommand{\qtwodryingtimespread}{26\,s}
\newcommand{\qtwogridrefinementdiff}{$3.2\times10^{-7}$\,kg/kg}
\newcommand{\qtwomassbalancerel}{$4.37\times10^{-13}$}
\newcommand{\qtwosensDplusfivepct}{-4.2031\,\%}
\newcommand{\qtwosensTinfplustwoC}{-6.0776\,\%}
\newcommand{\qtwosenshmplusfivepct}{-0.5332\,\%}
\newcommand{\qtwosensmassequationform}{-14.7245\,\%}
\newcommand{\qtwotableregionenvinsensitivity}{$2.64\times10^{-5}$\,kg/kg}
\newcommand{\qtwotimerefinementdiff}{$2.78\times10^{-4}$\,h}
\newcommand{\qthreeCcenterfivefourh}{0.1539\,kg/kg}
\newcommand{\qthreeCcenteratend}{0.15\,kg/kg}
\newcommand{\qthreeCsurfaceatend}{0.0527\,kg/kg}
\newcommand{\qthreeconsistencywithQtwo}{57.4647\,h}
\newcommand{\qthreedryingtime}{57.4647\,h}
\newcommand{\qthreegridrefinementdiff}{$3.2\times10^{-7}$\,kg/kg}
\newcommand{\qthreemassbalancerel}{$4.76\times10^{-13}$}
\newcommand{\qthreetimerefinementdiff}{$3.6\times10^{-3}$\,h}
\newcommand{\qfourRatdryingendcm}{1.2\,cm}
\newcommand{\qfourbaselinedryingtimeseconds}{15\,s}
\newcommand{\qfourbaselinetablesixdiff}{$4\times10^{-4}$\,kg/kg}
\newcommand{\qfourdryingtimehours}{52.6361\,h}
\newcommand{\qfourgridrefinementdiff}{$2.12\times10^{-6}$\,kg/kg}
\newcommand{\qfourmassbalancerel}{$2.56\times10^{-14}$}
\newcommand{\qfourmovingboundarydev}{$1.42\times10^{-13}$\,kg/kg}
\newcommand{\qfourradiusinterphours}{$2.8\times10^{-3}$\,h}
\newcommand{\qfourreferencemaxabserr}{$1.15\times10^{-6}$\,kg/kg}
\newcommand{\qfoursensDplusfivepcthours}{-2.1722\,h}
\newcommand{\qfoursensTinfplustwoChours}{-3.1306\,h}
\newcommand{\qfoursenshmplusfivepcthours}{-0.1278\,h}
\newcommand{\qfoursensmassequationformhours}{-9.225\,h}
\newcommand{\qfourshrinkagepercent}{40.1\,\%}
\newcommand{\qfourtablesixendcenter}{0.15\,kg/kg}
\newcommand{\qfourtablesixendsurface}{0.0526\,kg/kg}
\newcommand{\qfourtablesixtsixcenter}{1.9473\,kg/kg}
\newcommand{\qfourtablesixtsixsurface}{0.5654\,kg/kg}
\newcommand{\qfourtimerefinementhours}{0\,h}
\newcommand{\qfourvsqthreedryingtimehours}{-4.8286\,h}

\title{中药材热风烘干过程的温湿耦合建模与尺寸收缩效应分析}
\author{}
\date{}

\linespread{1.4}\selectfont   % 正文行距，对齐获奖论文的 20pt \u884c间距
\captionsetup{font={small},labelsep=quad,justification=centering}
\captionsetup[table]{position=below}
\begin{document}
\maketitle

\begin{abstract}
本文研究圆柱形中药材在热风烘干过程中温度与水分浓度的时空演化规律，并确定烘干所需时长。
针对题目给出的四个问题，本文建立\textbf{统一的一维径向轴对称耦合模型}，
采用守恒型的顶点有限体积格式离散，并用四类相互独立的证据验证全部数值结果。

\textbf{针对问题一}，建立预热平衡阶段的温湿耦合模型，以附件 1 的实测环境序列驱动边界。
结果表明该阶段\textbf{以升温为主、失水仅限表层}：$1800$~s 时中心与表面温度分别为 
\qoneTcenteroneeightzerozeros 与 \qoneTsurfaceoneeightzerozeros，
而中心干基含水率仍保持初值 \qoneCcenteroneeightzerozeros，仅表面降至 \qoneCsurfaceoneeightzerozeros。
其物理原因在于水分扩散的特征时间约为 $22$~h，远长于本阶段的三十分钟。

\textbf{针对问题二和问题三}，改用温度与浓度相关的物性经验式，建立含双向耦合的全过程模型，
并以三步后向 Euler 启动的 Crank--Nicolson 格式推进。结果表明烘干过程由\textbf{内部扩散控制}，
全域最大干基含水率于 \qtwodryingtime 首次降至判据以下，即烘干时长为 \qtwodryingtime；
问题三独立求解所得 \qthreedryingtime 与之完全一致，构成一次跨问自检。

\textbf{针对问题四}，药材在失水过程中会发生尺寸收缩。本文引入归一化径向坐标，
把移动边界转化为已知的表观对流项，从而在固定网格上保持精确的守恒性。
结果表明烘干时长为 \qfourdryingtimehours，此时半径已收缩至 \qfourRatdryingendcm（收缩 \qfourshrinkagepercent）。
\textbf{这一结论揭示了一个反直觉的机理}：附录 4 给出的扩散系数前因子比附录 3 小 $5.7$ 倍，
若只替换物性本应\textbf{延长}烘干时间；但收缩把扩散路径缩短到约 $0.6$ 倍，其效应压倒了扩散系数的下降，
最终使烘干时长比不考虑收缩的情形缩短 \qfourvsqthreedryingtimehours。
\textbf{因此尺寸收缩不可忽略，问题四的烘干时长不能由问题三直接推算。}

\textbf{模型检验}方面，本文不依赖先验论证，而给出四类证据：与常物性圆柱扩散的解析级数解对照，
最大绝对误差达 \qfourreferencemaxabserr 量级；网格与时间的双向加密使报告值不再变化；
另写一套空间离散与时间推进都不同的格心实现独立复算，报告点差异在一个末位单位以内；
并把离散质量守恒写成可逐时步核验的恒等式，残差达 \qfourmassbalancerel。
对问题四，本文进一步设计\textbf{解析不变量检验}——关闭表面传质后收缩不得改变干基含水率，
实测偏离仅 \qfourmovingboundarydev，使移动边界项的正确性可以被证伪。

\textbf{灵敏度分析}给出了一个对结论表述至关重要的结果：烘干时长对题目给定的表面传质系数并不敏感
（放大 $5\%$ 仅变化 \qtwosenshmplusfivepct），而主要由模型口径决定
（质量方程读法变化 \qtwosensmassequationform，环境温度升高 $2^\circ$C 变化 \qtwosensTinfplustwoC）。
因此本文的全部结论均以条件式给出，并显式列出恒温段环境为外推口径、端面轴向输运被忽略等局限。

\noindent\textbf{关键词}：热风烘干；径向轴对称；有限体积法；移动边界；归一化坐标；参数灵敏度
\end{abstract}

\newpage

% Sections are \input fragments produced by paper-section-writer
\section{问题重述}

\subsection{问题背景}

中药材的干燥是决定成品品质的关键工序之一，热风烘干因设备简单、易于控制而成为常用方式。
整个过程分为预热平衡与恒温干燥两个阶段：预热阶段药材由室温逐步升温、内部水分开始向表面迁移；
恒温干燥阶段烘房维持稳定的温湿环境，水分持续从表面蒸发、内部水分不断向外补充，直至达到干燥要求。
工艺参数选取不当会带来干燥效率低、能耗高、成品品质不稳定等问题，而传统的逐批试验优化成本高、周期长，
因此借助机理分析与数值仿真获得药材内部温度与水分浓度的时空规律，具有重要的工程意义。

在实际干燥过程中，药材会因水分流失而收缩，水分扩散路径随之缩短；与此同时，扩散系数本身又受温度和浓度影响。二者对干燥时长的影响方向相反、相互竞争。因此，只有通过建模进行定量分析，才能判断哪一效应占主导。

\subsection{问题提出}

题目围绕圆柱形中药材在热风烘干过程中的温度场与水分浓度场提出四个问题：

\begin{itemize}[label=\smash{\large\ding{226}}, leftmargin=2.4em, itemsep=3pt]
  \item \textbf{问题一}：给出烘干初期 $0\sim1800$~s 内药材内部温度与水分浓度的变化规律，
        并按题目规定的时刻与位置列表。同时输出该时段内每隔固定时间、固定径向间隔的完整计算结果。
  \item \textbf{问题二}：在恒温环境下建立整个烘干过程中温湿场变化的模型，
        按题目规定的时刻与位置给出结果，并输出完整计算结果。
  \item \textbf{问题三}：依据题目给定的干燥要求确定恒温干燥所需的时间，
        并按题目规定的时刻与位置给出水分浓度。
  \item \textbf{问题四}：在药材随失水而收缩的条件下确定烘干时长，
        并按题目规定的时刻与位置给出水分浓度，输出完整计算结果。
\end{itemize}

\section{问题的分析}

\subsection{总体分析}

四个问题的物理对象相同，都是圆柱形药材内部的热量与水分传递，因此可以建立\textbf{统一的一维径向模型骨架}，
再按问题在三个维度上分支：\textbf{物性口径}（问题一用附录 2、问题二与问题三用附录 3、问题四用附录 4）、
\textbf{时程与环境}（问题一是 $30$~min 的预热段、问题二与问题三是跨越多日的恒温干燥段）、
以及问题四新增的\textbf{几何口径}（计算域随失水而收缩）。这样安排可以避免同一套方程被重复推导四遍，
也使四问的结果具有直接可比性。总体技术路线见图~\ref{fig:chain}。

\begin{figure}[H]
  \centering
  \includegraphics[width=0.97\textwidth]{fig06_method_and_roles.pdf}
  \caption{四个问题的建模主线与角色分工。自下而上为：统一的控制方程骨架，
  四问在物性口径与时程上的分支，每一问的主方法、可用基线与独立验证（三者来自不同的空间离散与
  时间推进，彼此不共享数值输入），以及贯穿四问的验证闭环。}
  \label{fig:chain}
\end{figure}

\subsection{各问的关键难点}

\textbf{问题一}的难点不在方程本身，而在\textbf{表面值的精度}：预热段仅 $30$~min，水分几乎只在极薄的表层响应，
表面位置承担了全部的离散误差预算；若把最外层网格中心的值直接当作表面值，会独自吃掉整个误差预算。
此外题面要求的 $100$~s 并非附件的采样时刻，必须插值。

\textbf{问题二与问题三}的难点有三：一是物性随温度与浓度变化，温度与水分\textbf{双向耦合}，
传热与传质不能再解耦求解；二是时程长达数日，显式时间推进的稳定步长在长时程下不可接受；
三是干燥判据位于含水率曲线的\textbf{平坦尾部}，判据穿越时刻对数值噪声高度敏感，必须单独评估其不确定性。

\textbf{问题四}的难点是\textbf{移动边界}：计算域随半径缩小而内移，若仍在固定网格上求解，
会在边界处引入人工的质量源或汇，使烘干时长出现系统性偏差。

\subsection{技术路线与验证策略}

本文的处理思路是：先建立归一化的守恒离散格式，使\textbf{质量守恒可以在离散层面被逐时步核验}；
再针对移动边界引入归一化坐标，把未知的移动边界转化为已知的表观对流项；
最后用四类相互独立的证据确认数值结果，而不依赖先验论证——解析解对照、网格与时间加密、
独立编写的第二套实现交叉验证、以及解析不变量检验。四问的验证结果汇总见图~\ref{fig:verify}。

\begin{figure}[H]
  \centering
  \includegraphics[width=0.92\textwidth]{fig05_numerical_verification.pdf}
  \caption{数值可信度的四类证据。(a) 网格加密残差与解析解对照残差；
  (b) 问题四的冻结半径算例与 Bessel 级数解析解的对照；
  (c) 主方法与独立实现的逐点最大差，虚线为一个末位单位对应的容差；
  (d) 问题四关闭表面传质后干基含水率的偏移量与其阈值。}
  \label{fig:verify}
\end{figure}

\section{模型假设}

本文的建模建立在下列假设之上。
每一条假设的影响都在第 6 节中量化。

\begin{enumerate}[label=\arabic*., leftmargin=2.6em, itemsep=4pt]
  \item \textbf{局部热力学平衡}：药材内部各点满足局部平衡，烘房环境的温度与水分浓度可以作为
        边界驱动量；否则第三类边界条件无法建立。
  \item \textbf{一维径向轴对称}：药材为长径比大于 $12$ 的圆柱，热风沿轴向与周向均匀作用，
        故只保留径向坐标，忽略两端面的轴向输运。该简化在 $30$~min 的预热窗口内影响不足长度的 $2\%$，
        但在数日的干燥时程内轴向渗透可达数厘米每端，其影响在检验部分量化。
  \item \textbf{质量方程取标准的 Fick 形式}：按题目明文的写法定式，密度、比热容与导热系数只进入能量方程。
        若改用与湿基密度严格自洽的写法，等效扩散系数会相差一个随浓度变化的因子，其对结果的影响在检验部分量化。
  \item \textbf{物性与经验式取题目给定值}：密度、比热容、导热系数、扩散系数与对流传热传质系数
        均按题目附录取值，不作拟合。
  \item \textbf{恒温干燥段的环境为定值}：附件 1 只给到 $14400$~s，其后按已确认的口径
        维持恒定的环境温度与水分浓度；该设定由附件 1 稳定段的实测均值核实。
  \item \textbf{忽略水分汽化潜热}：题目未给潜热数据，能量方程中不含相变项，
        该简化会使升温阶段的预测偏快。
  \item \textbf{变物性系数取上一时间层}：使每个时间步仍是线性方程组，代价是时间格式降为一阶，
        其精度由时间加密实验核验。
  \item \textbf{收缩各向同性且严格按附件 2}：问题四的半径随时间的变化完全由附件 2 决定，
        收缩只改变几何尺寸，不改变干基含水率的定义。
  \item \textbf{问题四的报告坐标取当前物理距离}：按题目已确认的口径，
        每个时刻按当时半径换算物理坐标，超出当前半径的位置不给出数值。
  \item \textbf{问题四的表观对流系数由附件 2 数值微分得到}：其噪声放大风险由\textbf{解析不变量检验}控制。
\end{enumerate}

\section{符号说明}

本文使用的符号与单位统一列于表~\ref{tab:symbols}。全文温度用摄氏度、含水率用\textbf{干基}
（每千克干物质所含的水分质量），长度在内部计算中为米、在结果表中为厘米，时间在交付文件中为秒、在表格中为小时。
为节省篇幅，表~\ref{tab:symbols} 只列出后文反复使用的符号，一次性符号在首次出现处就地说明。

\begin{longtable}{@{}l p{7.3cm} l@{}}
  \caption{符号说明}\label{tab:symbols}\\
  \toprule
  符号 & 含义 & 单位 \\
  \midrule
  \endfirsthead
  \toprule 符号 & 含义 & 单位 \\ \midrule \endhead
  \bottomrule \endfoot
  $R_0,\ R(t),\ R_{\min}$ & 药材的初始半径、当前半径、收缩终值 & cm \\
  $L$ & 药材长度 & cm \\
  $r$ & 到药材中心的径向距离（物理坐标） & cm \\
  $\xi$ & 归一化径向坐标，等于当前物理距离与当前半径之比 & --- \\
  $t$ & 时间 & s（交付文件）/ h（表格） \\
  $T(r,t)$ & 药材内部温度 & $^\circ$C \\
  $C(r,t)$ & 药材内部干基水分浓度 & kg/kg \\
  $T_0,\ C_0$ & 温度与含水率的初值 & $^\circ$C, kg/kg \\
  $T_\infty(t),\ C_\infty(t)$ & 烘房环境温度与环境水分浓度 & $^\circ$C, kg/kg \\
  $\rho$ & 药材密度 & kg/m$^3$ \\
  $c_p$ & 药材比热容 & J/(kg$\cdot$K) \\
  $k$ & 药材导热系数 & W/(m$\cdot$K) \\
  $\alpha$ & 热扩散率，等于导热系数除以密度与比热容之积 & m$^2$/s \\
  $D$ & 水分扩散系数 & m$^2$/s \\
  $h$ & 对流传热系数 & W/(m$^2\cdot$K) \\
  $h_m$ & 对流传质系数 & m/s \\
  $\dot R$ & 药材半径的变化率 & m/s \\
  $N$ & 空间网格数（主方法取 $800$） & --- \\
  $\Delta r,\ \Delta\xi$ & 物理坐标与归一化坐标下的网格间距 & m, --- \\
  $\Delta t$ & 时间步长 & s \\
  $w_j$ & 第 $j$ 个对偶控制体的归一化度量，等于该控制体上径向坐标的积分 & --- \\
  $k_c$ & 表观对流系数，等于半径变化率与半径之比 & s$^{-1}$ \\
  $W$ & 单位长度药材内的总水量 & kg/m \\
  $C_{\text{target}}$ & 烘干判据，即全域最大干基含水率的上限 & kg/kg \\
  $t_{\text{dry}}$ & 烘干时长 & h \\
  $\text{Bi},\ \text{Bi}_m,\ \text{Fo}$ & 传热、传质 Biot 数与 Fourier 数 & --- \\
\end{longtable}

\section{模型的建立与数值方法}

\subsection{建模准备与总体框架}

四个问题研究的是同一个对象，即圆柱形药材在热风中的升温与失水，控制方程、边界条件与初值完全一致。
因此我们不逐问重复建模，而是先建立一套统一的径向轴对称模型，四个问题都在它上面求解，
只在物性经验式、环境时程与计算区域三处分支：问题一取问题给定的常物性，
问题二与问题三取附录 3 的温度—浓度双向耦合经验式，问题四再叠加由附件 2 决定的尺寸收缩。
\\这样处理既免去重复推导，也使四问的结果具有直接可比性。本节先给出控制方程与定解条件，
再给出统一的数值方法，各问的具体参数与求解结果分别见第六章至第九章。

\subsection{控制方程的建立}

\subsubsection{由守恒律导出控制方程}

药材内部既无内热源也无化学反应，故热量只靠导热传递，水分只由浓度梯度驱动。
由微元法，在药材内取半径由 $r$ 到 $r+\Delta r$ 的一层薄环，并沿轴向取单位长度，由\textbf{能量守恒定律}环内热量的变化率应等于
两侧界面净导入的导热热流，即

\begin{equation}
\rho c_p\frac{\partial T}{\partial t}\cdot 2\pi r\,\Delta r
=\left.\left(k\cdot 2\pi r\frac{\partial T}{\partial r}\right)\right|_{r+\Delta r}
-\left.\left(k\cdot 2\pi r\frac{\partial T}{\partial r}\right)\right|_{r}.
\label{eq:energybalance}
\end{equation}

两端同除以 $2\pi r\,\Delta r$，再令 $\Delta r\to 0$，右端即成为散度形式，于是得到能量方程

\begin{equation}
\rho c_p\frac{\partial T}{\partial t}
=\frac{1}{r}\frac{\partial}{\partial r}\left(k\,r\,\frac{\partial T}{\partial r}\right).
\label{eq:energy}
\end{equation}

对水分做同样的守恒分析：环内水量的变化率等于两侧界面净导入的水分质量流，
而质量流的本构关系即\textbf{ Fick 定律}，即通量等于扩散系数乘以浓度梯度的负值，于是

\begin{equation}
\frac{\partial C}{\partial t}
=\frac{1}{r}\frac{\partial}{\partial r}\left(D\,r\,\frac{\partial C}{\partial r}\right).
\label{eq:mass}
\end{equation}

观察到，问题一中物性 $k$、$D$ 与状态无关，式~\eqref{eq:energy} 与式~\eqref{eq:mass} 彼此独立，
温度场与水分场可以分别求解。

需要说明的是，式~\eqref{eq:mass} 中不含密度：本文按题目明文的 Fick 标准形式建立质量方程，
密度、比热容与导热系数只出现在能量方程中。这一读法决定了后文全部时长结论的适用范围，
其影响量级在灵敏度分析中定量给出。

\subsubsection{边界条件与初值}

\begin{tabular}{@{}p{1.4em}p{0.85\textwidth}@{}}
① & 在中心 \(r=0\) 处，由轴对称性给出零梯度条件；\\[2pt]
② & 在表面 \(r=R\) 处，热流与质量流均由对流控制。\\
\end{tabular}

把外法向通量写开，即得第三类（Robin）边界条件
\begin{equation}
\left.-k\frac{\partial T}{\partial r}\right|_{r=R}
=h\left(T(R,t)-T_\infty(t)\right),
\label{eq:bc}
\end{equation}

\begin{equation}
\left.-D\frac{\partial C}{\partial r}\right|_{r=R}
=h_m\left(C(R,t)-C_\infty(t)\right).
\label{eq:bcm}
\end{equation}

其中 $h$ 与 $h_m$ 分别为对流传热系数与对流传质系数，$T_\infty$ 与 $C_\infty$ 为烘房环境的
温度与水分浓度。初始时刻温度与含水率在全域均匀，即 $T(r,0)=T_0$、$C(r,0)=C_0$。\\
[2pt]环境序列由附件 1 的实测数据给出（$0$--$14400$~s，步长 $60$~s），采样点之间的时刻按相邻两点的分段线性插值取值。
$14400$~s 之后记录终止，本文取末段（约 $t\ge9000$~s）稳定平台的实测均值并保持不变，
即 $T_\infty=50$~$^\circ$C、$C_\infty=0.05$~kg/kg，以支撑问题二至问题四 的计算。

\subsubsection{特征时间与量级分析}

数值求解之前，先由控制方程判断两个阶段的主控机制，为时间步长与验证设计定标。
对式~\eqref{eq:mass} 与式~\eqref{eq:bcm} 作无量纲化，取 $\xi=r/R$、$u=(C-C_\infty)/(C_0-C_\infty)$，得

\begin{equation}
\frac{\partial u}{\partial\tau}
=\frac{1}{\xi}\frac{\partial}{\partial\xi}\left(\xi\frac{\partial u}{\partial\xi}\right),\qquad
\tau=\frac{Dt}{R^2},
\label{eq:massnd}
\end{equation}

\begin{equation}
\left.-\frac{\partial u}{\partial\xi}\right|_{\xi=1}
=\text{Bi}_m\,u(1,\tau),\qquad
\text{Bi}_m=\frac{h_mR}{D}=\frac{R/D}{1/h_m}.
\label{eq:bcnd}
\end{equation}

式~\eqref{eq:massnd} 不含任何参数，物性只由式~\eqref{eq:bcnd} 进入，
可见水分过程仅由 Fourier 数 $\tau$ 与传质 Biot 数 $\text{Bi}_m$ 支配。

$\tau$ 表示已过去的时间相当于多少个特征扩散时间 $R^2/D$。预热阶段处于湿端
（$C\approx1.5$--$2.5$~kg/kg、$T\approx28$~$^\circ$C），$R^2/D\approx20$--$22$~h，
而该阶段只有 $30$~min，$t=1800$~s 时 $\tau\approx0.025$、穿透深度 $\sqrt{Dt}\approx0.3$~cm，
不足半径的两成，故\textbf{失水只能发生在表层}。
\\$\text{Bi}_m$ 是内部扩散阻力与外部对流传质阻力之比，本题湿端约 $3$、干端约 $20$，全程大于 $1$，
说明\textbf{表面传质阻力小于内部扩散阻力，干燥由内部扩散控制}。

这两个量级判断将在第六、七章的灵敏度分析中被定量证实。

\subsection{数学求解过程}

四个问题共用同一条求解链条：把空间导数离散成守恒通量形式，把时间导数按 $\theta$ 加权格式推进，
问题四再叠加移动边界的表观对流项，每个时间步归结为一个三对角方程组并用追赶法求解，
最后在题目要求的报告时刻与位置取值。四问的差别只在代入哪套物性经验式、
是否启用坐标变换以及计算区域如何取，链条本身不变。完整流程见图~\ref{fig:flow}。

\begin{figure}[H]
  \centering
  \includegraphics[width=0.94\textwidth]{fig07_solution_flow.pdf}
  \caption{统一求解流程。由环境序列与几何数据出发，经控制体权重与界面系数构造，
  按问题选取物性经验式与时间格式，最后由追赶法逐时步推进并输出报告值。}
  \label{fig:flow}
\end{figure}

下面按空间离散、时间离散、移动边界处理与验证设计的顺序给出各环节。

\subsubsection{空间离散：对偶控制体的顶点有限体积格式}

空间离散采用有限体积法：把守恒律写在每个小控制体上，界面通量一进一出相互抵消，
因此水分不会凭空产生或消失。下面按“取控制体—写通量”两步构造格式。

把区间 $[0,R]$ 等分为 $N$ 段，节点为 $r_j=j\Delta r$。以每个节点为中心取对偶控制体：
内部节点的控制体为 $[r_j-\Delta r/2,\;r_j+\Delta r/2]$，首尾节点各取半格。
把式~\eqref{eq:mass} 在每个控制体上作径向加权积分（权为 $r$，即环面积因子），得

\begin{equation}
\int_{V_j} r\frac{\partial C}{\partial t}\,dr
=\left[\,r D\frac{\partial C}{\partial r}\right]_{\text{两侧界面}},
\label{eq:fv}
\end{equation}

右端只出现界面上的通量，这正是有限体积法的出发点。把左端的积分近似为
$w_j\,\mathrm{d}C_j/\mathrm{d}t$，其中权重 $w_j=\int_{V_j}r\,\mathrm{d}r$ 为对偶控制体上径向坐标的积分，
把右端的界面通量用相邻节点的差商近似，即得半离散方程

\begin{equation}
w_j\frac{\mathrm{d}C_j}{\mathrm{d}t}
=r_{j+\frac12}D_{j+\frac12}\frac{C_{j+1}-C_j}{\Delta r}
-r_{j-\frac12}D_{j-\frac12}\frac{C_j-C_{j-1}}{\Delta r}.
\label{eq:semi}
\end{equation}

图~\ref{fig:mesh} 给出这一离散的几何含义：左侧在圆柱截面上取一行径向节点，
右侧放大后可以看到节点、对偶控制体与界面通量的相对位置。
本文取 $R=2.0$~cm、$N=800$ 个节点，故 $\Delta r=R/N=25$~$\mu$m。

\begin{figure}[H]
  \centering
  \includegraphics[width=0.98\textwidth]{fig00_mesh_schematic.pdf}
  \caption{顶点型有限体积离散示意}
  \label{fig:mesh}
\end{figure}

\subsubsection{两个端点的处理}

在中心 $r=0$ 处，对偶控制体的内侧界面半径为 $0$，式~\eqref{eq:semi} 中对应的通量项自动为零，
轴对称条件\textbf{由格式本身满足}，无需对 $1/r$ 的奇点作任何特殊处理。
在表面 $r=R$ 处，节点本身就是边界点，把式~\eqref{eq:bcm} 的通量代入式~\eqref{eq:fv} 的界面项，得

\begin{equation}
w_N\frac{\mathrm{d}C_N}{\mathrm{d}t}
=r_{N-\frac12}D_{N-\frac12}\frac{C_{N-1}-C_N}{\Delta r}
-R\,h_m\left(C_N-C_\infty\right).
\label{eq:surf}
\end{equation}

式~\eqref{eq:surf} 说明：报告的表面值就是解变量 $C_N$ 本身，不需要由最外层单元中心回算。
这避免了半单元重构引入的一阶误差，也正是我们选择顶点格式作为主要求解格式的原因。

\subsubsection{时间离散：$\theta$ 加权与启动步}

对半离散系统 $\mathbf{w}\,\mathrm{d}\mathbf{y}/\mathrm{d}t=\mathbf{Q}(\mathbf{y})$ 取 $\theta$ 加权，得

\begin{equation}
\mathbf{y}^{n+1}-\mathbf{y}^{n}
=\Delta t\left[\theta\,\mathbf{Q}^{n+1}+(1-\theta)\mathbf{Q}^{n}\right].
\label{eq:theta}
\end{equation}

当 $\theta=1$ 时为后向 Euler，精度一阶但耗散强、无条件稳定；当 $\theta=1/2$ 时为
Crank--Nicolson，精度二阶且 A 稳定。本文取 $\theta=1/2$，但\textbf{前 $3$ 步改用 $\theta=1$}：
初值处的温湿场与边界驱动之间存在突变，纯 Crank--Nicolson 会把这一突变转化为随时间衰减的振荡，
以后向 Euler 启动即可将其抹平；三步之后解已充分光滑，再切回二阶格式不影响整体精度。

由于物性依赖温度与含水率，每一步的系数取上一时间层的值，使每个时间步仍是一个形如

\begin{equation}
a_jy_{j-1}+b_jy_j+c_jy_{j+1}=d_j
\label{eq:tridiag}
\end{equation}

的\textbf{三对角方程组}，可用追赶法在 $O(N)$ 时间内求解。系数滞后带来的时间一阶误差
由时间加密实验核验（见第七、八章）。

\subsubsection{移动边界的坐标变换}

问题四中半径随时间收缩为 $R(t)$，固定网格无法贴合移动的边界。引入归一化坐标 $\xi=r/R(t)$，
把未知的移动域 $[0,R(t)]$ 映射为固定的 $[0,1]$，这一步把“网格跟着边界走”转化为“边界不动、
方程多一项”。设 $C(r,t)=\tilde C(\xi,t)$，由链式法则

\begin{equation}
\left.\frac{\partial C}{\partial t}\right|_r
=\left.\frac{\partial \tilde C}{\partial t}\right|_\xi
+\frac{\partial \tilde C}{\partial \xi}\cdot\left.\frac{\partial \xi}{\partial t}\right|_r,
\qquad \left.\frac{\partial \xi}{\partial t}\right|_r=-\frac{r\dot R}{R^2}=-\frac{\xi\dot R}{R},
\label{eq:chain}
\end{equation}

同时 $\partial/\partial r=R^{-1}\,\partial/\partial\xi$，代入式~\eqref{eq:mass} 即得

\begin{equation}
\left.\frac{\partial \tilde C}{\partial t}\right|_\xi
=\frac{1}{\xi R^2}\frac{\partial}{\partial \xi}\left(D\,\xi\,\frac{\partial \tilde C}{\partial \xi}\right)
+\underbrace{\frac{\xi\dot R}{R}\frac{\partial \tilde C}{\partial \xi}}_{\text{表观对流项}} .
\label{eq:moving}
\end{equation}

这一项的物理含义是清楚的：干基含水率按定义不随几何收缩而改变，因此在固定于材料的归一化坐标下，
收缩表现为一个把浓度剖面按 $\xi\dot R/R$ 平流的对流项，而不是一个凭空产生或消灭水分的源汇。
$\dot R$ 由附件 2 的单调插值解析求导得到，是已知函数，因此式~\eqref{eq:moving} 仍是线性方程组，
无需重构网格即可求解。

\subsubsection{移动边界的守恒性核验与不变量检验}

移动边界下的守恒性不易目测，本文把它写成一个可以直接核验的离散恒等式。
把式~\eqref{eq:moving} 两端同乘 $\xi$ 并在 $[0,1]$ 上积分，左端为
$\frac{\mathrm{d}}{\mathrm{d}t}\int_0^1\xi\tilde C\,\mathrm{d}\xi$（$\xi$ 域固定，可移入导数）。
右端第一项按分部积分给出界面项 $\left[\xi D\,\partial_\xi\tilde C/R^2\right]_0^1
=-\frac{h_m}{R}(C_N-C_\infty)$，第二项给出
$\frac{\dot R}{R}\int_0^1\xi^2\partial_\xi\tilde C\,\mathrm{d}\xi
=\frac{\dot R}{R}\left(C_N-2\int_0^1\xi\tilde C\,\mathrm{d}\xi\right)$。
把两个积分用归一化坐标下对偶控制体的权重 $\omega_j=\int_{\text{cell}}\xi\,\mathrm{d}\xi$ 离散
（它与式~\eqref{eq:fv} 中的 $w_j$ 相差常数因子 $R^2$），整理得

\begin{equation}
\frac{\mathrm{d}}{\mathrm{d}t}\sum_j \omega_jC_j
=-2\frac{\dot R}{R}\sum_j \omega_jC_j+\frac{\dot R}{R}C_N-\frac{h_m}{R}\left(C_N-C_\infty\right).
\label{eq:conservation}
\end{equation}

式~\eqref{eq:conservation} 是一个\textbf{可以直接逐时步核验的离散恒等式}：
左端由数值解给出，右端由当步的系数给出，二者之差即格式的守恒性残差。

再设一条更锋利的判据：令 $h_m=0$ 但保持收缩，则式~\eqref{eq:moving} 退化为纯平流方程，
其解析解要求 $\tilde C(\xi,t)\equiv C_0$ 恒成立——因为干基含水率不随几何缩小而改变。
若表观对流项的符号或系数写错，收缩就会凭空产生或消灭水分，该检验必然立刻暴露。
这里我们把这条判据称为\textbf{解析不变量检验}，它是针对移动边界项最有针对性的一条验证设计，
其结果见第九章。

\subsubsection{求解流程与统一验证设计}

综合以上推导，四个问题共用同一套求解流程：读取环境序列与几何数据；构造对偶控制体权重与界面系数；
按问题选择物性经验式与时间格式；组装三对角方程组并用追赶法逐时步推进；在报告时刻按题目要求的
位置取值输出。其中问题四额外按 $\dot R$ 叠加表观对流项。完整流程见图~\ref{fig:flow}。

每一问都配齐三条相互独立的实现路径：本文方法（顶点有限体积）、对照方法（格心有限体积，
其表面值由式~\eqref{eq:bcm} 回算）与独立复算（常物性下的级数解析解，以及式~\eqref{eq:conservation}
给出的守恒恒等式与本节的不变量检验）。三条路径不共享数值输入，只在其结果上相互比对，
以保证“算对了”这一判断不是自证。各问的比对结果分别见第六至第九章。

\section{模型的检验}

本节从误差来源、收敛性、参数与口径灵敏度、以及不变量四个方面检验前节的数值结果。
检验所依据的全部数值均由各问的运行记录直接给出，未作任何人工调整。

\subsection{误差来源分析}

数值解的误差可以分解为四项：空间离散误差、时间离散误差、解析近似误差（物性经验式本身的不确定性）、
以及与独立实现比对时暴露的实现差异。前三项通过加密实验与解析对照评估，第四项通过第二套实现交叉验证评估。
需要强调的是，本题的\textbf{主要不确定性并不来自数值误差}，而来自题目给定的物性与环境口径——
这一点在灵敏度分析中被定量地证实。

\subsection{收敛性检验}

对每一问分别做空间与时间的独立加密：空间把网格数由 $400$ 加密到 $800$，时间把内部步长减半或把求解器容差收紧一个量级。
加密后题目要求的位置上的四位小数不再变化，说明所报数值已达到网格与步长无关的水平。
此外，在常物性、固定半径的极限下，本文的实现与圆柱扩散问题的 Bessel 级数解析解对照，
四问的最大绝对误差分别达到 
\qoneanalyticalmaxabserr、\qtwoanalyticalmaxabserr（问题三与其同模型）、\qtwoanalyticalmaxabserr（问题三同模型） 与 \qfourreferencemaxabserr。

需要说明的是，问题一使用的是自适应时间积分，不存在固定步长，因此其时间精度的证据是\textbf{求解器容差加密}
（收紧一个量级后报告点上的表值变化为 \qonetimeaccuracytolerancediff）；
第 (2)(3)(4) 问使用固定步长，证据是步长减半后\textbf{烘干时长}的变化，分别为 
\qtwotimerefinementdiff、\qthreetimerefinementdiff 与 \qfourtimerefinementhours。
问题四另需检验半径插值口径的影响：把附件 2 的单调插值换成线性插值，烘干时长仅变化 \qfourradiusinterphours。

\subsection{参数与口径灵敏度分析}

灵敏度分析针对两类对象：一类是题目给定但数值本身带不确定性的参数（对流传热与传质系数、扩散系数），
另一类是建模过程中必须做出选择的\textbf{口径}（环境外推方式、质量方程的读法、半径插值方式）。
在恒温干燥的基准工况下，各单项扰动对烘干时长的影响汇总于表~\ref{tab:sens}。

\begin{table}[H]
  \centering
  \caption{单项扰动对烘干时长的影响（基准工况的烘干时长为 \qtwodryingtime）}
  \label{tab:sens}
  \small
  \begin{tabular}{lccc}
    \toprule
    扰动项 & 问题二与问题三 & 问题四 & 结论 \\
    \midrule
    对流传质系数放大 $5\%$ & \qtwosenshmplusfivepct & \qfoursenshmplusfivepcthours & 影响很小 \\
    扩散系数放大 $5\%$ & \qtwosensDplusfivepct & \qfoursensDplusfivepcthours & 中等影响 \\
    环境温度升高 $2^\circ$C & \qtwosensTinfplustwoC & \qfoursensTinfplustwoChours & 较大影响 \\
    \bottomrule
  \end{tabular}
\end{table}

表~\ref{tab:sens} 给出的是一个反直觉但对结论至关重要的结果：\textbf{烘干时长对表面传质系数并不敏感}
（在 $5\%$ 的扰动下只变化半个百分点），原因是本题的干燥由\textbf{内部水分扩散控制}，
传质 Biot 数约为 $20$，表面传质阻力不是瓶颈；
因此本文的所有时长结论都写成条件式：\textbf{在题目给定的经验式与质量方程读法之下}。
问题一的规律不同：预热段的表面值对传质系数敏感，放大 $5\%$ 即使表面含水率变化 
\qonehmplusfivepctsurfaceshift，故该问的表面值只在题目给定值下成立。

\subsection{不变量与守恒性检验}

离散格式的守恒性不依赖网格加密，可以\textbf{逐时步核验}。本文把质量守恒写成一个离散恒等式，
四问的残差分别为 \qonemassbalancerel、\qtwomassbalancerel、\qthreemassbalancerel 与 \qfourmassbalancerel，
均达到机器精度量级。

对问题四的移动边界，本文另设一条\textbf{解析不变量检验}：关闭表面传质、只让计算域收缩，
干基含水率必须严格保持不变。若表观对流项的符号或系数写错，收缩就会凭空产生或消灭水分，
该检验因此可以直接\textbf{证伪}项本身的正确性。实测的最大偏离仅为 \qfourmovingboundarydev。

\subsection{与独立实现的交叉验证}

每一问都用一套空间离散与时间推进都不同的实现独立复算：基线采用格心有限体积，表面值由对流边界关系回算，
问题二与问题三的基线还不使用启动步。两套实现在题目要求的报告点上的最大差异为 
\qoneindependentimplementationsdiff（问题一）、\qtwodryingtimespread（问题二与问题三的烘干时长）、
以及 \qfourbaselinetablesixdiff 与 \qfourbaselinedryingtimeseconds（问题四的表值与时长）。

问题二与问题三的时长差异值得单独说明：两套实现在报告点上一致到四位小数的舍入边界，
却仍带来数十秒的烘干时刻差。这不是实现缺陷，而是\textbf{判据位于含水率曲线平坦尾部}的必然结果——
曲线在判据附近近乎水平，极小的数值差异即可换来可观的时间差。因此本文在给出烘干时长的同时，
始终给出该量级的不确定度。

\subsection{跨问一致性}

问题二与问题三的模型、参数与判据完全相同，因此两问应当给出同一个烘干时长。
实测两问分别为 \qtwodryingtime 与 \qthreedryingtime，完全一致（问题三的对照量为 \qthreeconsistencywithQtwo），
这构成一次\textbf{跨问自检}。问题三判据穿越前的最后一个采样点中心含水率为 \qthreeCcenterfivefourh，
穿越后中心与表面分别为 \qthreeCcenteratend 与 \qthreeCsurfaceatend，说明穿越是一次性的、不存在反复越界。

\subsection{检验结论}

四类证据一致表明：本文报告的数值已达到所要求的精度，且其不确定性的主要来源已被识别并量化。
表~\ref{tab:verify} 汇总了四问的全部检验量，图~\ref{fig:verify} 给出四类证据的图形汇总。

\begin{table}[H]
  \centering
  \caption{四问的检验量汇总}\label{tab:verify}
  \footnotesize
  \begin{tabularx}{\textwidth}{l *{5}{>{\centering\arraybackslash}X}}
    \toprule
    问题 & 解析对照/(kg/kg) & 网格加密/(kg/kg) & 时间加密 & 独立实现差 & 离散守恒 \\
    \midrule
    第 (1) 问 & \qoneanalyticalmaxabserr & \qonegridrefinementdiff & \qonetimeaccuracytolerancediff & \qoneindependentimplementationsdiff & \qonemassbalancerel \\
    第 (2) 问 & \qtwoanalyticalmaxabserr & \qtwogridrefinementdiff & \qtwotimerefinementdiff & \qtwodryingtimespread & \qtwomassbalancerel \\
    第 (3) 问 & \qthreedryingtime & \qthreegridrefinementdiff & \qthreetimerefinementdiff & \qthreeconsistencywithQtwo & \qthreemassbalancerel \\
    第 (4) 问 & \qfourreferencemaxabserr & \qfourgridrefinementdiff & \qfourtimerefinementhours & \qfourbaselinedryingtimeseconds & \qfourmassbalancerel \\
    \bottomrule
  \end{tabularx}
  \par\smallskip\footnotesize 第 (3) 问的解析对照取与其同模型的第 (2) 问结果；
  第 (4) 问的时间加密一列为内部步长减半引起的烘干时长变化。
\end{table}


\begin{figure}[H]
  \centering
  \includegraphics[width=0.92\textwidth]{fig05_numerical_verification.pdf}
  \caption{数值可信度的四类证据。(a) 网格加密残差与解析解对照残差；
  (b) 问题四的冻结半径算例与 Bessel 级数解析解的对照；
  (c) 主方法与独立实现的逐点最大差，虚线为一个末位单位对应的容差；
  (d) 问题四关闭表面传质后干基含水率的偏移量与其阈值。}
  \label{fig:verify}
\end{figure}
\section{模型的评价与推广}

\subsection{模型的优点}

\textbf{其一，四问共用一套可核验的守恒离散框架。}本文没有为每个问题单独编写求解器，
而是把控制方程骨架、边界处理与守恒格式统一起来，只在物性口径、时程与坐标口径上分支。
这样既避免了同一物理过程被重复推导，也使四问的结果具有直接可比性。

\textbf{其二，质量守恒可以在离散层面被逐时步核验。}顶点有限体积的通量形式使全域水量的变化
与表面通量之间存在精确的离散恒等式，残差达到机器精度；对移动边界问题，本文进一步给出
“关闭传质后收缩不改变干基含水率”的解析不变量检验，使表观对流项的正确性可以被\textbf{证伪}，
而不是仅凭收敛曲线判断。

\textbf{其三，主导不确定性被识别并量化。}分析表明烘干时长对题目给定的传质系数并不敏感，
而主要由模型口径决定。这一结论改变了结论的呈现方式：论文不应给出一个孤立的烘干时长数字，
而应同时给出它所依赖的口径与相应的影响幅度。

\textbf{其四，报告的数值可逐条溯源。}所有进入论文的数值都来自运行记录，并由机器生成的宏引用，
没有一个是手工抄录的；表格亦由脚本从运行结果直接生成，从根本上排除了抄写误差。

\subsection{模型的缺点与改进方向}

\textbf{其一，维度简化。}模型忽略了药材两端面的轴向输运。在 $30$~min 的预热窗口内这在数量级上成立，
但在数日的干燥时程内轴向渗透可达数厘米每端，必须作为局限声明。若要改进，可把模型扩展为二维轴对称问题，
或在端面处补充轴向通量的等效边界条件。

\textbf{其二，质量方程的读法。}本文按题目明文的 Fick 标准形式建式，密度不进入质量方程，
因此几何收缩在模型内不是质量汇。替代读法会使烘干时长变化近 $15\%$，该差异只能通过口径声明而非数值手段消除。

\textbf{其三，经验式的适用性。}物性完全采用题目给定的经验式，未作独立验证；
扩散系数的强非线性使结果对经验式的形式较为敏感，这是模型本身无法回避的不确定性。

\textbf{其四，判据尾部的平坦性。}干燥判据位于含水率曲线的平坦段，使烘干时长对数值噪声敏感。
改进方向是改用\textbf{时间区间}而非单个时刻来报告烘干完成，或同时给出判据附近的曲线形态。

\subsection{模型的推广}

本文的建模框架不限于中药材。凡是满足“圆柱（或球形）几何、内部扩散控制、表面为对流边界”的干燥、
吸附或浸提过程，都可以直接套用：把物性经验式替换为目标材料的数据，把环境序列替换为实际的工艺曲线，
并保留本文的守恒检验与不变量检验即可。问题四的归一化坐标方法尤其具有推广价值——
凡计算域随时间单调变化（收缩、膨胀、溶解、烧蚀）且变化规律已知的扩散问题，都可以用同一手法
把移动边界转化为已知的表观对流项，从而在固定网格上求解并保留精确的守恒性。


% AI-use declaration (injected from decision ledgers / ai_use_disclosure.md)


\begin{thebibliography}{99}
\bibitem{virtanen2020scipy} Virtanen, Pauli and Gommers, Ralf and Oliphant, Travis E. and others. SciPy 1.0: fundamental algorithms for scientific computing in Python. Nature Methods. 2020
\bibitem{harris2020numpy} Harris, Charles R. and Millman, K. Jarrod and van der Walt, St{\'e}fan and others. Array programming with NumPy. Nature. 2020
\bibitem{carslaw1959conduction} Carslaw, Horatio S. and Jaeger, John C.. Conduction of Heat in Solids. Oxford University Press. 1959
\bibitem{incropera2007heat} Incropera, Frank P. and DeWitt, David P. and Bergman, Theodore L. and Lavine, Adrienne S.. Fundamentals of Heat and Mass Transfer. John Wiley \& Sons. 2007
\bibitem{leveque2002finite} LeVeque, Randall J.. Finite Volume Methods for Hyperbolic Problems. Cambridge University Press. 2002
\bibitem{crank1947practical} Crank, John. The Mathematics of Diffusion. Oxford University Press. 1975
\bibitem{patankar1980numerical} Patankar, Suhas V.. Numerical Heat Transfer and Fluid Flow. Hemisphere Publishing. 1980
\end{thebibliography}

\clearpage
\appendix
\section{支撑材料与程序代码}

\subsection{支撑材料清单}

\begin{table}[H]
  \centering
  \caption{程序文件清单}\label{tab:code}
  \small
  \begin{tabularx}{\textwidth}{@{}llX@{}}
    \toprule
    文件 & 对应问题 & 作用 \\
    \midrule
    \texttt{code/Q1/q1\_common.py} & 问题一 & 公共常数、输入读取、报告点映射与解析级数解 \\
    \texttt{code/Q1/q1\_main.py} & 问题一 & 主方法：顶点有限体积 + 自适应隐式积分；产出结果表与完整场文件 \\
    \texttt{code/Q1/q1\_baseline.py} & 问题一 & 可用基线：格心有限体积 + 固定步长 Crank--Nicolson \\
    \texttt{code/Q1/q1\_verifier.py} & 问题一 & 独立验证：解析级数参照与不变量核验 \\
    \texttt{code/Q2/q2\_main.py} 等 & 问题二 & 含双向耦合的全过程求解、基线与验证 \\
    \texttt{code/Q3/q3\_main.py} 等 & 问题三 & 烘干时长的确定、基线与验证 \\
    \texttt{code/Q4/q4\_main.py} & 问题四 & 主方法：归一化坐标下的移动边界求解，含表观对流项 \\
    \texttt{code/Q4/q4\_baseline.py} & 问题四 & 可用基线：归一化坐标下的格心实现 \\
    \texttt{code/Q4/q4\_verifier.py} & 问题四 & 独立验证：自建级数解与报告口径逐格核验 \\
    \texttt{code/figures/} & 全部 & 论文图表的生成脚本（含统一的视觉系统与渲染自检） \\
    \bottomrule
  \end{tabularx}
\end{table}

四个问题的完整计算文件为 \texttt{result1.xlsx}--\texttt{result4.xlsx}，
分别给出题目要求的时间步长与径向间隔下的温度场与水分浓度场。

下文给出核心程序全文。为便于对照正文，每段代码的注释均标出其在论文中的对应位置
（如"论文对应：5.2 节"）。问题二与问题三复用问题一、问题四的求解器，其差异仅在物性
经验式与环境条件，故不再重复列出，完整工程见支撑材料目录 \texttt{code/}。

\subsection{问题一公共模块}
\lstinputlisting[style=pycode, caption={问题一公共模块 \texttt{q1\_common.py}：常数、附件读取、报告点映射与解析级数参照}]
  {code/Q1/q1_common.py}

\subsection{问题一主方法}
\lstinputlisting[style=pycode, caption={问题一主方法 \texttt{q1\_main.py}：顶点有限体积离散与自适应 BDF 时间推进}]
  {code/Q1/q1_main.py}

\subsection{问题一独立验证}
\lstinputlisting[style=pycode, caption={问题一独立验证 \texttt{q1\_verifier.py}：自建 Bessel 参照与契约不变量核验}]
  {code/Q1/q1_verifier.py}

\subsection{问题四主方法}
\lstinputlisting[style=pycode, caption={问题四主方法 \texttt{q4\_main.py}：\texorpdfstring{$\xi=r/R(t)$}{xi = r/R(t)} 归一化坐标下的移动边界有限体积求解}]
  {code/Q4/q4_main.py}

\subsection{问题四独立验证}
\lstinputlisting[style=pycode, caption={问题四独立验证 \texttt{q4\_verifier.py}：自建级数参照、交付文件逐格核验}]
  {code/Q4/q4_verifier.py}

\end{document}
